%! Matrix Computations and A = LU — Unit 2, Lesson 2.3 — Guided Notes (Answer Key)
\documentclass[10pt]{article}
\usepackage{linalg-article}
\usepackage{linalg-key}   % pulls in -boxes; do NOT also load -boxes
\newcommand{\TallMath}[1]{$\displaystyle #1\rule[-1.4em]{0pt}{3.2em}$}
\newcommand{\xx}{\mathbf{x}}
\newcommand{\bb}{\mathbf{b}}
\newcommand{\cc}{\mathbf{c}}

\begin{document}

\pageheader{Unit 2, Lesson 2.3}{Guided Notes --- Answer Key}

\vspace{0.12in}

\begin{objectivebox}
\textbf{By the end of this lesson, I will be able to\dots}
\begin{itemize}\setlength{\itemsep}{2pt}
  \item factor a matrix as $A=LU$ --- $U$ is elimination's result, $L$ \emph{stores the multipliers};
  \item \emph{check} a factorization by multiplying $LU$ back to $A$;
  \item solve $A\xx=\bb$ in two sweeps: forward-solve $L\cc=\bb$, then back-solve $U\xx=\cc$;
  \item explain why $A=LU$ is the efficient way to re-solve one system for many right-hand sides.
\end{itemize}
\end{objectivebox}

\vspace{0.10in}

\begin{vocabbox}
\small
Fill in each term as we build it together.
\par\vspace{2pt}
\vocabans{Lower triangular $L$}{$1$s on the diagonal, the multipliers $\ell$ below, $0$s above; entry $\ell_{ij}$ is the multiplier that cleared position $(i,j)$}
\vocabans{Upper triangular $U$}{the result of forward elimination: pivots on the diagonal, $0$s below}
\vocabans{Factorization $A=LU$}{splitting $A$ into (lower)$\times$(upper); $L$ remembers how we eliminated, $U$ remembers what we got}
\vocabans{Forward substitution}{solve a lower-triangular system $L\cc=\bb$ from the top row down}
\vocabans{Two triangular sweeps}{solve $A\xx=\bb$ as $L\cc=\bb$ (forward) then $U\xx=\cc$ (back) --- no inverse}
\end{vocabbox}

\begin{hookbox}
Same bakery as last lesson: the recipe matrix $A$ never changes --- only the order $\bb$ does.
\[
  A = \begin{bmatrix} 1 & 3 \\ 2 & 7 \end{bmatrix}, \qquad
  \text{this week } \bb = \begin{bmatrix} 7 \\ 15 \end{bmatrix}.
\]
\begin{itemize}\setlength{\itemsep}{1pt}
  \item In 2.2 we built the whole inverse $A^{-1}$. But elimination in 2.1 already did the hard part.
        What if we could just \emph{save} that work and reuse it for every order?
        \ansline{Save elimination as $A=LU$: keep the result $U$ and the multipliers in $L$, then reuse them on any new $\bb$.}
\end{itemize}
\end{hookbox}

\begin{notesbox}{1. Saving Elimination as $A=LU$}
Run elimination on $A=\begin{bsmallmatrix} 1&3\\2&7 \end{bsmallmatrix}$ once (pivot $1$, multiplier
$\ell=2$). Row~2 becomes (row~2)$-2\,$(row~1), leaving the upper-triangular
\[
  U = \begin{bmatrix} 1 & 3 \\ 0 & 1 \end{bmatrix}.
\]
Now build the lower-triangular $\mathbf{L}$: start from the identity, put $1$s on the diagonal, and
drop the multiplier $\ell=2$ into the spot it cleared --- \emph{with a plus sign} (this is the
undo-matrix $E^{-1}$ from 2.2):
\[
  L = \begin{bmatrix} 1 & 0 \\ {\color{keyred}\mathbf{2}} & 1 \end{bmatrix}.
\]
\textbf{Check the factorization} by multiplying $LU$ --- it should rebuild $A$:
\[
  LU = \begin{bmatrix} 1 & 0 \\ 2 & 1 \end{bmatrix}\begin{bmatrix} 1 & 3 \\ 0 & 1 \end{bmatrix}
  = \begin{bmatrix} {\color{keyred}\mathbf{1}} & {\color{keyred}\mathbf{3}}\\[2pt]{\color{keyred}\mathbf{2}} & {\color{keyred}\mathbf{7}}\end{bmatrix} = A. \quad\checkmark
\]
So $A=LU$: $U$ remembers \emph{what elimination got}; $L$ remembers \emph{how we did it}.
\end{notesbox}

\begin{notesbox}{2. $L$ Just Stores the Multipliers (a $3\times3$)}
Here is the payoff: building $L$ takes \textbf{no extra work}. Whatever multipliers you use during
elimination, you simply drop each one into $L$ (with $1$s on the diagonal). Eliminate
$A=\begin{bsmallmatrix} 1&2&1\\2&5&5\\3&7&8 \end{bsmallmatrix}$, recording each multiplier
$\ell=\dfrac{\text{entry}}{\text{pivot}}$:
\[
  \begin{array}{ll}
    \ell_{21} = 2 \div 1 = 2, & \text{(clears the $(2,1)$ entry)}\\[3pt]
    \ell_{31} = 3 \div 1 = {\color{keyred}\mathbf{3}}, & \text{(clears the $(3,1)$ entry)}\\[3pt]
    \ell_{32} = 1 \div 1 = {\color{keyred}\mathbf{1}}. & \text{(clears the $(3,2)$ entry)}
  \end{array}
\]
Elimination leaves $U=\begin{bsmallmatrix} 1&2&1\\0&1&3\\0&0&2 \end{bsmallmatrix}$ (pivots $1,1,2$). Now
place each multiplier into its spot in $L$:
\[
  L = \begin{bmatrix} 1 & 0 & 0 \\ \ell_{21} & 1 & 0 \\ \ell_{31} & \ell_{32} & 1 \end{bmatrix}
    = \begin{bmatrix} 1 & 0 & 0 \\ {\color{keyred}\mathbf{2}} & 1 & 0 \\ {\color{keyred}\mathbf{3}} & {\color{keyred}\mathbf{1}} & 1 \end{bmatrix}.
\]
The multipliers slot straight in --- $L$ is a free record of the elimination.
\end{notesbox}

\begin{notesbox}{3. Solving $A\xx=\bb$ in Two Triangular Sweeps}
Once $A=LU$, we \emph{never} touch $A$ again. Since $A\xx=\bb$ means $L(U\xx)=\bb$, give the inside
piece a name: $\cc=U\xx$. Then solve two easy triangular systems back-to-back.

\vspace{3pt}
\textbf{Sweep 1 --- forward, $L\cc=\bb$} (top row down). For the bakery, $\bb=\begin{bsmallmatrix} 7\\15 \end{bsmallmatrix}$:
\[
  \begin{bmatrix} 1 & 0 \\ 2 & 1 \end{bmatrix}\begin{bmatrix} c_1 \\ c_2 \end{bmatrix} = \begin{bmatrix} 7 \\ 15 \end{bmatrix}
  \;\Rightarrow\; c_1 = {\color{keyred}\mathbf{7}},\quad 2c_1 + c_2 = 15 \Rightarrow c_2 = {\color{keyred}\mathbf{1}}.
\]
So $\cc=\begin{bsmallmatrix} 7\\1 \end{bsmallmatrix}$ --- exactly the updated right-hand side elimination produced in 2.2!
The forward sweep \emph{is} ``update $\bb$ as you eliminate.''

\vspace{3pt}
\textbf{Sweep 2 --- back, $U\xx=\cc$} (bottom row up):
\[
  \begin{bmatrix} 1 & 3 \\ 0 & 1 \end{bmatrix}\begin{bmatrix} x \\ y \end{bmatrix} = \begin{bmatrix} 7 \\ 1 \end{bmatrix}
  \;\Rightarrow\; y = {\color{keyred}\mathbf{1}},\quad x + 3y = 7 \Rightarrow x = {\color{keyred}\mathbf{4}}.
\]
So $(x,y)=(\,\ans{$4,\,1$}\,)$ --- the same answer 2.2 got, with no inverse in sight.
\end{notesbox}

\boxguard
\begin{notesbox}{4. Why Factor? The ``Matrix Computations'' Point}
Factoring $A=LU$ is the \emph{expensive} step --- but you do it \textbf{once}. After that, every new
order $\bb$ costs only the two quick sweeps above. Compare the options for the same $A$ with many
right-hand sides:
\begin{itemize}\setlength{\itemsep}{2pt}
  \item Re-run elimination from scratch each time --- repeats all the hard work every order.
  \item Build $A^{-1}$ (Lesson 2.2) --- powerful, but more work to build than $L$ and $U$.
  \item \textbf{Factor once, then two sweeps per order} --- cheapest for solving. \emph{This is what
        computers actually do.}
\end{itemize}
$L$ and $U$ are just elimination, \emph{saved} and ready to reuse.
\end{notesbox}

\boxguard
\begin{practicebox}
Let $A=\begin{bmatrix} 1 & 4 \\ 3 & 10 \end{bmatrix}$ and $\bb=\begin{bmatrix} 6 \\ 16 \end{bmatrix}$.
\begin{enumerate}[leftmargin=1.6em]
  \item Find the multiplier, then write $L$ and $U$. \hfill $\ell=\ans{$3$}$
\begin{work}
  U &= \begin{bsmallmatrix} 1&4\\0&-2 \end{bsmallmatrix} \\
  L &= \begin{bsmallmatrix} 1&0\\3&1 \end{bsmallmatrix}
\end{work}
  \item \textbf{Check} $LU=A$.
\begin{work}
  \begin{bsmallmatrix} 1&0\\3&1 \end{bsmallmatrix}\begin{bsmallmatrix} 1&4\\0&-2 \end{bsmallmatrix}
      &= \begin{bsmallmatrix} 1&4\\3&10 \end{bsmallmatrix} = A \ \ \text{\checkmark}
\end{work}
  \item Solve $A\xx=\bb$ by two sweeps ($L\cc=\bb$, then $U\xx=\cc$). \hfill $x=\ans{$2$}$\quad $y=\ans{$1$}$
\begin{work}
  c_1 &= 6 \\
  3(6)+c_2 &= 16 \\
  c_2 &= -2 \\
  -2y &= -2 \\
  y &= 1 \\
  x+4 &= 6,\ x = 2
\end{work}
\end{enumerate}
\end{practicebox}

\end{document}
